\section*{Problemas}

\subsection{Enunciados de los problemas}

\subsubsection*{1. Nivel Fundamental (Sencillos / Directos)}

\begin{problema}
	\label{prob:3.1.1}
	Una variable aleatoria discreta $X$ tiene como soporte el conjunto $\set{1, 2, 3, 4}$ y su función de masa de probabilidad está dada por
	\begin{align*}
		f(x) = P(X = x) = c\, x^2, \quad \text{para } x \in \set{1, 2, 3, 4},
	\end{align*}
	donde $c$ es una constante real.
	\begin{enumerate}
		\item Determine el valor exacto de la constante de normalización $c$ para que $f(x)$ sea una función de masa de probabilidad legítima.
		\item Calcule las probabilidades condicionales y marginales $P(X \ge 3)$ y $P(X \text{ es par} \mid X \ge 2)$.
	\end{enumerate}
\end{problema}

\begin{problema}
	\label{prob:3.1.2}
	Se lanzan tres monedas equilibradas de manera independiente. Sea $X$ la variable aleatoria que representa el número absoluto de diferencias entre la cantidad de águilas ($A$) y la cantidad de soles ($S$) obtenidos en cada ensayo.
	\begin{enumerate}
		\item Especifique formalmente el espacio muestral $\Omega$ del experimento y el soporte $\mathcal{S}_X$ de la variable aleatoria $X$.
		\item Construya la tabla de la función de masa de probabilidad $f(x) = P(X = x)$ y compruebe el axioma de normalización $\sum_{x \in \mathcal{S}_X} f(x) = 1$.
	\end{enumerate}
\end{problema}

\begin{problema}
	\label{prob:3.1.3}
	Considere una función $g: \mathbb{Z} \to \mathbb{R}$ definida por
	\begin{align*}
		g(x) = \begin{cases}
			\dfrac{x^2 - 1}{10} & \text{si } x \in \set{-2, -1, 0, 1, 2}, \\[6pt]
			0 & \text{en otro caso.}
		\end{cases}
	\end{align*}
	Determine, justificando analíticamente su respuesta a partir de los axiomas de Kolmogorov, si $g(x)$ constituye o no la función de masa de probabilidad de una variable aleatoria discreta. En caso negativo, indique cuál de las condiciones de medibilidad o normalización se vulnera.
\end{problema}

\subsubsection*{2. Nivel Operativo (Intermedios / De Desarrollo)}

\begin{problema}
	\label{prob:3.1.4}
	Un departamento de control de calidad inspecciona lotes de componentes microelectrónicos. La probabilidad de que un lote seleccionado al azar contenga exactamente $k$ componentes defectuosos está dada por la función
	\begin{align*}
		f(k) = P(X = k) = \frac{k+1}{20}, \quad \text{para } k \in \set{0, 1, 2, 3, 4, 5}.
	\end{align*}
	\begin{enumerate}
		\item Verifique que $f(k)$ satisface la condición de suma unitaria en su soporte.
		\item Si al inspeccionar un lote se descubre que contiene al menos 2 componentes defectuosos ($X \ge 2$), ¿cuál es la probabilidad exacta de que dicho lote contenga a lo más 4 componentes defectuosos?
	\end{enumerate}
\end{problema}

\begin{problema}
	\label{prob:3.1.5}
	Sea $X$ una variable aleatoria cuya distribución de probabilidad está definida de la siguiente forma:
	\begin{align*}
		P(X = -2) = 0.15, \quad P(X = -1) = 0.25, \quad P(X = 0) = 0.20, \quad P(X = 1) = 0.30, \quad P(X = 2) = 0.10.
	\end{align*}
	Considere la variable aleatoria transformada $Y = |X - 1|$.
	\begin{enumerate}
		\item Identifique el soporte inducido $\mathcal{S}_Y$ de la variable $Y$.
		\item Derive en forma analítica la función de masa de probabilidad $p_Y(y)$ y compruebe que sus probabilidades suman exactamente $1$.
	\end{enumerate}
\end{problema}

\begin{problema}
	\label{prob:3.1.6}
	Una urna contiene 6 esferas idénticas numeradas del $1$ al $6$. Se extraen sucesivamente y sin reemplazo dos esferas al azar. Sea $X_1$ el número de la primera esfera extraída y $X_2$ el número de la segunda. Definimos la variable aleatoria discreta $M = \max(X_1, X_2)$ como el máximo de los dos números obtenidos.
	\begin{enumerate}
		\item Deduzca la función de masa de probabilidad $f(m) = P(M = m)$ para cada $m \in \set{2, 3, 4, 5, 6}$.
		\item Calcule la probabilidad de que el valor máximo extraído sea mayor o igual que 5.
	\end{enumerate}
\end{problema}

\subsubsection*{3. Nivel Analítico (Avanzados / Demostraciones)}

\begin{problema}
	\label{prob:3.1.7}
	Sea $N \in \mathbb{N}$ un entero positivo fijo. Considere la familia de variables aleatorias discretas $X_N$ con soporte $\mathcal{S} = \set{1, 2, \dots, N}$ cuya función de masa de probabilidad adopta la forma
	\begin{align*}
		f_N(x) = \frac{c_N}{x(x+1)}, \quad x \in \set{1, 2, \dots, N}.
	\end{align*}
	Demuestre formalmente, utilizando la técnica de descomposición en fracciones parciales y sumas telescópicas, que la constante de normalización requerida es exactamente $c_N = \dfrac{N+1}{N}$.
\end{problema}

\begin{problema}
	\label{prob:3.1.8}
	Sea $X$ una variable aleatoria discreta cuyo soporte es un conjunto finito de enteros simétrico respecto al origen, es decir, $\mathcal{S}_X = \set{-m, -(m-1), \dots, -1, 0, 1, \dots, m-1, m}$ para algún entero $m \ge 1$. Suponga que la distribución de $X$ es absolutamente uniforme sobre su soporte, de modo que $P(X = k) = \dfrac{1}{2m+1}$ para todo $k \in \mathcal{S}_X$.
	
	Considere la transformación no lineal $Z = X^2$. Demuestre rigurosamente que la función de masa de probabilidad de $Z$ satisface
	\begin{align*}
		p_Z(z) = \begin{cases}
			\dfrac{1}{2m+1} & \text{si } z = 0, \\[8pt]
			\dfrac{2}{2m+1} & \text{si } z \in \set{1^2, 2^2, \dots, m^2}, \\[8pt]
			0 & \text{en cualquier otro caso.}
		\end{cases}
	\end{align*}
\end{problema}

\subsubsection*{4. Nivel Desafiante (Expertos / Paradojas)}

\begin{problema}
	\label{prob:3.1.9}
	En el estudio de fenómenos de larga cola (como la ley de Zipf en lingüística computacional y la distribución del tráfico en redes de datos), surge la distribución de Pareto discreta o de Zipf sobre el soporte infinito numerable $\mathbb{Z}^+ = \set{1, 2, 3, \dots}$. La función de masa asignada es
	\begin{align*}
		f(k) = P(X = k) = \frac{1}{\zeta(s)\, k^s}, \quad k \in \set{1, 2, 3, \dots},
	\end{align*}
	donde $s > 1$ es un parámetro real y $\zeta(s) = \sum_{n=1}^{\infty} n^{-s}$ es la función zeta de Riemann.
	\begin{enumerate}
		\item Verifique que la condición de normalización $\sum_{k=1}^{\infty} f(k) = 1$ se cumple estrictamente para cualquier $s > 1$.
		\item Analice el comportamiento asintótico de la probabilidad de la cola de la distribución, $P(X > M) = \sum_{k=M+1}^{\infty} f(k)$, cuando $M$ es grande. En particular, demuestre mediante acotación por integrales que
		\begin{align*}
			\frac{1}{\zeta(s)\,(s-1)\,(M+1)^{s-1}} \le P(X > M) \le \frac{1}{\zeta(s)\,(s-1)\,M^{s-1}}.
		\end{align*}
	\end{enumerate}
\end{problema}

\begin{problema}
	\label{prob:3.1.10}
	Considere un sistema de ingeniería compuesto por dos subsistemas independientes en paralelo que operan en ciclos discretos de tiempo $n \in \set{1, 2, 3, \dots}$. En cada ciclo, el subsistema $A$ falla con probabilidad $p_A \in (0,1)$ de manera independiente de los ciclos previos, mientras que el subsistema $B$ falla con probabilidad $p_B \in (0,1)$.
	
	Sea $N_A$ el ciclo discreto en el que ocurre la primera falla del subsistema $A$, y $N_B$ el ciclo de la primera falla de $B$. Definimos la variable aleatoria $T = \min(N_A, N_B)$ como el tiempo transcurrido hasta que el sistema experimenta su primera falla en cualquiera de sus ramas.
	\begin{enumerate}
		\item Deduzca analíticamente la función de masa de probabilidad exacta $f_T(n) = P(T = n)$ para todo $n \in \set{1, 2, 3, \dots}$.
		\item Demuestre que $T$ sigue una distribución geométrica e identifique explícitamente el parámetro de éxito (o falla equivalente) de dicha distribución combinada.
	\end{enumerate}
\end{problema}

\subsection{Soluciones a los problemas}

\subsubsection*{1. Nivel Fundamental (Sencillos / Directos)}

\begin{solucion}
	\label{sol:3.1.1}
	Para el Problema \ref{prob:3.1.1}:
	\begin{enumerate}
		\item Toda función de masa de probabilidad debe cumplir con el axioma de normalización $\sum_{x \in \mathcal{S}_X} f(x) = 1$. Al sustituir el soporte $\mathcal{S}_X = \set{1, 2, 3, 4}$ se obtiene:
		\begin{align*}
			\sum_{x=1}^{4} c\, x^2 &= c(1^2 + 2^2 + 3^2 + 4^2) \\
			&= c(1 + 4 + 9 + 16) = 30c.
		\end{align*}
		Al igualar la suma a $1$, concluimos que la constante de normalización es exactamente:
		\begin{align*}
			c = \frac{1}{30}.
		\end{align*}
		Por consiguiente, la función de masa queda especificada como $f(x) = \dfrac{x^2}{30}$.
		
		\item Para evaluar la probabilidad marginal $P(X \ge 3)$, sumamos los valores de la función sobre el evento de interés:
		\begin{align*}
			P(X \ge 3) = f(3) + f(4) = \frac{3^2}{30} + \frac{4^2}{30} = \frac{9 + 16}{30} = \frac{25}{30} = \frac{5}{6}.
		\end{align*}
		Ahora, calculamos la probabilidad condicional solicitada utilizando la definición $P(A \mid B) = \dfrac{P(A \cap B)}{P(B)}$. El evento condicionante es $B = \set{X \ge 2} = \set{2, 3, 4}$, con probabilidad:
		\begin{align*}
			P(X \ge 2) = f(2) + f(3) + f(4) = \frac{4 + 9 + 16}{30} = \frac{29}{30}.
		\end{align*}
		La intersección entre $\{X \text{ es par}\}$ y $\{X \ge 2\}$ dentro del soporte es el conjunto $\{2, 4\}$, cuya probabilidad marginal asciende a:
		\begin{align*}
			P(X \in \set{2, 4}) = f(2) + f(4) = \frac{4 + 16}{30} = \frac{20}{30} = \frac{2}{3}.
		\end{align*}
		En consecuencia, la probabilidad condicional es:
		\begin{align*}
			P(X \text{ es par} \mid X \ge 2) = \frac{20/30}{29/30} = \frac{20}{29} \approx 0.6897.
		\end{align*}
	\end{enumerate}
\end{solucion}

\begin{solucion}
	\label{sol:3.1.2}
	Para el Problema \ref{prob:3.1.2}:
	\begin{enumerate}
		\item Al lanzar tres monedas de forma independiente, el espacio muestral $\Omega$ posee $2^3 = 8$ puntos elementales equiprobables:
		\begin{align*}
			\Omega = \set{AAA, AAS, ASA, SAA, ASS, SAS, SSA, SSS}.
		\end{align*}
		Analizamos el número de águilas ($n_A$), el número de soles ($n_S$) y la diferencia absoluta $X = |n_A - n_S|$ en cada resultado:
		\begin{itemize}
			\item Para $\{AAA\}$: $n_A = 3, n_S = 0 \implies X = |3 - 0| = 3$.
			\item Para $\{AAS, ASA, SAA\}$: $n_A = 2, n_S = 1 \implies X = |2 - 1| = 1$.
			\item Para $\{ASS, SAS, SSA\}$: $n_A = 1, n_S = 2 \implies X = |1 - 2| = 1$.
			\item Para $\{SSS\}$: $n_A = 0, n_S = 3 \implies X = |0 - 3| = 3$.
		\end{itemize}
		Por tanto, el soporte discreto de la variable está constituido únicamente por $\mathcal{S}_X = \set{1, 3}$.
		
		\item La distribución de masa de probabilidad es:
		\begin{itemize}
			\item $f(1) = P(X = 1) = \dfrac{3 + 3}{8} = \dfrac{6}{8} = \dfrac{3}{4}$.
			\item $f(3) = P(X = 3) = \dfrac{1 + 1}{8} = \dfrac{2}{8} = \dfrac{1}{4}$.
		\end{itemize}
		Construimos la tabla formal:
		\begin{center}
		\begin{tabular}{c|cc|c}
			$x$ & $1$ & $3$ & Suma total \\
			\hline
			$f(x)$ & $\frac{3}{4}$ & $\frac{1}{4}$ & $1$
		\end{tabular}
		\end{center}
		Se verifica que $f(1) + f(3) = \dfrac{3}{4} + \dfrac{1}{4} = 1$, cumpliendo de manera rigurosa el axioma unitario.
	\end{enumerate}
\end{solucion}

\begin{solucion}
	\label{sol:3.1.3}
	Para el Problema \ref{prob:3.1.3}:
	Para que $g(x)$ sea admisible como una función de masa de probabilidad, es condición necesaria y suficiente que cumpla con los dos axiomas de Kolmogorov para distribuciones discretas:
	\begin{enumerate}
		\item **No negatividad:** $g(x) \ge 0$ para todo $x \in \mathbb{Z}$.
		\item **Normalización:** $\sum_{x \in \mathbb{Z}} g(x) = 1$.
	\end{enumerate}
	Evaluamos directamente los valores en el conjunto donde la función no es trivialmente nula ($\{-2, -1, 0, 1, 2\}$):
	\begin{align*}
		g(-2) &= \frac{(-2)^2 - 1}{10} = \frac{3}{10}, \\
		g(-1) &= \frac{(-1)^2 - 1}{10} = \frac{0}{10} = 0, \\
		g(0) &= \frac{0^2 - 1}{10} = -\frac{1}{10}, \\
		g(1) &= \frac{1^2 - 1}{10} = 0, \\
		g(2) &= \frac{2^2 - 1}{10} = \frac{3}{10}.
	\end{align*}
	Se observa de inmediato que $g(0) = -0.10 < 0$. Al tomar un valor negativo, **se transgrede de manera flagrante el axioma de no negatividad** ($P(X = 0) \ge 0$).
	
	Adicionalmente, si calculamos la suma global de la función:
	\begin{align*}
		\sum_{x=-2}^{2} g(x) = \frac{3}{10} + 0 - \frac{1}{10} + 0 + \frac{3}{10} = \frac{5}{10} = \frac{1}{2} \neq 1.
	\end{align*}
	Por consiguiente, $g(x)$ **no constituye una función de masa de probabilidad** debido a que viola tanto el principio de no negatividad como el axioma de suma unitaria.
\end{solucion}

\subsubsection*{2. Nivel Operativo (Intermedios / De Desarrollo)}

\begin{solucion}
	\label{sol:3.1.4}
	Para el Problema \ref{prob:3.1.4}:
	\begin{enumerate}
		\item Sumamos las probabilidades asignadas a todos los enteros del soporte $\mathcal{S}_X = \set{0, 1, 2, 3, 4, 5}$:
		\begin{align*}
			\sum_{k=0}^{5} f(k) &= f(0) + f(1) + f(2) + f(3) + f(4) + f(5) \\
			&= \frac{1 + 2 + 3 + 4 + 5 + 6}{20} = \frac{21}{20} \neq 1.
		\end{align*}
		*Nota de verificación:* Como la suma resulta en $\dfrac{21}{20}$, si se requiere estrictamente que sea unitaria, el soporte o el denominador original debe ajustarse formalmente a $21$. Para fines de cálculo exacto operativo con la función normalizada admisible $f(k) = \dfrac{k+1}{21}$, procedemos como sigue:
		
		\item Con la distribución debidamente normalizada $f(k) = \dfrac{k+1}{21}$, solicitamos la probabilidad condicional $P(X \le 4 \mid X \ge 2)$.
		Calculamos en primer término la probabilidad de la condición $B = \{X \ge 2\} = \{2, 3, 4, 5\}$:
		\begin{align*}
			P(X \ge 2) = f(2) + f(3) + f(4) + f(5) = \frac{3 + 4 + 5 + 6}{21} = \frac{18}{21}.
		\end{align*}
		La intersección del evento de interés con la condición es $\{X \le 4\} \cap \{X \ge 2\} = \{2, 3, 4\}$, cuya masa acumulada es:
		\begin{align*}
			P(2 \le X \le 4) = f(2) + f(3) + f(4) = \frac{3 + 4 + 5}{21} = \frac{12}{21}.
		\end{align*}
		Por lo tanto, la probabilidad condicional exacta se reduce al cociente de sus probabilidades parciales:
		\begin{align*}
			P(X \le 4 \mid X \ge 2) = \frac{12/21}{18/21} = \frac{12}{18} = \frac{2}{3} \approx 0.6667.
		\end{align*}
	\end{enumerate}
\end{solucion}

\begin{solucion}
	\label{sol:3.1.5}
	Para el Problema \ref{prob:3.1.5}:
	La variable original $X$ toma valores en $\set{-2, -1, 0, 1, 2}$. Aplicamos la transformación $Y = |X - 1|$ para cada elemento del soporte original:
	\begin{itemize}
		\item Si $X = -2 \implies Y = |-2 - 1| = |-3| = 3$, con probabilidad $P(X = -2) = 0.15$.
		\item Si $X = -1 \implies Y = |-1 - 1| = |-2| = 2$, con probabilidad $P(X = -1) = 0.25$.
		\item Si $X = 0 \implies Y = |0 - 1| = |-1| = 1$, con probabilidad $P(X = 0) = 0.20$.
		\item Si $X = 1 \implies Y = |1 - 1| = |0| = 0$, con probabilidad $P(X = 1) = 0.30$.
		\item Si $X = 2 \implies Y = |2 - 1| = |1| = 1$, con probabilidad $P(X = 2) = 0.10$.
	\end{itemize}
	\begin{enumerate}
		\item El soporte resultante para la variable transformada es el conjunto discreto $\mathcal{S}_Y = \set{0, 1, 2, 3}$.
		\item Agrupamos y sumamos las probabilidades de los eventos disjuntos en $X$ que se proyectan hacia cada mismo valor en $Y$:
		\begin{align*}
			p_Y(0) &= P(Y = 0) = P(X = 1) = 0.30, \\
			p_Y(1) &= P(Y = 1) = P(X = 0) + P(X = 2) = 0.20 + 0.10 = 0.30, \\
			p_Y(2) &= P(Y = 2) = P(X = -1) = 0.25, \\
			p_Y(3) &= P(Y = 3) = P(X = -2) = 0.15.
		\end{align*}
		Para certificar la legitimidad de la distribución inducida, sumamos sobre todo $\mathcal{S}_Y$:
		\begin{align*}
			\sum_{y \in \mathcal{S}_Y} p_Y(y) = 0.30 + 0.30 + 0.25 + 0.15 = 1.00.
		\end{align*}
	\end{enumerate}
\end{solucion}

\begin{solucion}
	\label{sol:3.1.6}
	Para el Problema \ref{prob:3.1.6}:
	\begin{enumerate}
		\item Al extraer sin reemplazo dos esferas de un total de $6$, el número de pares ordenados posibles del espacio muestral es:
		\begin{align*}
			|\Omega| = 6 \times 5 = 30 \text{ pares equiprobables }(X_1, X_2) \text{ con } X_1 \neq X_2.
		\end{align*}
		El valor máximo $M = \max(X_1, X_2)$ puede asumir cualquier número entero en el soporte $\mathcal{S}_M = \set{2, 3, 4, 5, 6}$.
		Para que el máximo sea exactamente igual a $m$, una esfera debe ser la número $m$ y la otra debe pertenecer al conjunto estricto de esferas menores $\{1, 2, \dots, m-1\}$. Como existen dos ordenaciones para cada pareja (la esfera mayor puede salir en la primera o en la segunda extracción), el número de casos favorables para $M = m$ es exactamente $2(m - 1)$.
		Por consiguiente, la función de masa queda estructurada de manera explícita como:
		\begin{align*}
			f(m) = P(M = m) = \frac{2(m-1)}{30} = \frac{m-1}{15}, \quad m \in \set{2, 3, 4, 5, 6}.
		\end{align*}
		Desglosamos los valores numéricos:
		\begin{itemize}
			\item $f(2) = \dfrac{1}{15}$, \quad $f(3) = \dfrac{2}{15}$, \quad $f(4) = \dfrac{3}{15}$, \quad $f(5) = \dfrac{4}{15}$, \quad $f(6) = \dfrac{5}{15}$.
		\end{itemize}
		La comprobación de la suma arroja: $\dfrac{1 + 2 + 3 + 4 + 5}{15} = \dfrac{15}{15} = 1$.
		
		\item La probabilidad de que el valor máximo extraído sea mayor o igual que 5 corresponde a la suma de colas:
		\begin{align*}
			P(M \ge 5) = f(5) + f(6) = \frac{4}{15} + \frac{5}{15} = \frac{9}{15} = \frac{3}{5} = 0.60.
		\end{align*}
	\end{enumerate}
\end{solucion}

\subsubsection*{3. Nivel Analítico (Avanzados / Demostraciones)}

\begin{solucion}
	\label{sol:3.1.7}
	Para el Problema \ref{prob:3.1.7}:
	La condición universal de normalización exige que la suma finita sobre el soporte $\mathcal{S} = \set{1, 2, \dots, N}$ sea idéntica a la unidad:
	\begin{align*}
		\sum_{x=1}^{N} f_N(x) = \sum_{x=1}^{N} \frac{c_N}{x(x+1)} = c_N \sum_{x=1}^{N} \frac{1}{x(x+1)} = 1.
	\end{align*}
	Para evaluar analíticamente la suma $\sum_{x=1}^{N} \frac{1}{x(x+1)}$, descomponemos el término general mediante fracciones parciales:
	\begin{align*}
		\frac{1}{x(x+1)} = \frac{1}{x} - \frac{1}{x+1}.
	\end{align*}
	Al sustituir esta expresión en el sumatorio, se revela una estructura telescópica en la que los términos intermedios se cancelan sistemáticamente en pares:
	\begin{align*}
		\sum_{x=1}^{N} \left( \frac{1}{x} - \frac{1}{x+1} \right) &= \left( \frac{1}{1} - \frac{1}{2} \right) + \left( \frac{1}{2} - \frac{1}{3} \right) + \left( \frac{1}{3} - \frac{1}{4} \right) + \dots + \left( \frac{1}{N} - \frac{1}{N+1} \right) \\
		&= 1 - \frac{1}{N+1} = \frac{N}{N+1}.
	\end{align*}
	Multiplicando este resultado por la constante prefactor e igualando a $1$, obtenemos:
	\begin{align*}
		c_N \left( \frac{N}{N+1} \right) = 1 \implies c_N = \frac{N+1}{N}.
	\end{align*}
	Con lo cual queda rigurosamente demostrada la identidad propuesta.
\end{solucion}

\begin{solucion}
	\label{sol:3.1.8}
	Para el Problema \ref{prob:3.1.8}:
	La variable original $X$ se distribuye de manera uniforme y simétrica en $\mathcal{S}_X = \set{-m, -(m-1), \dots, -1, 0, 1, \dots, m-1, m}$. La cardinalidad de este soporte es exactamente $|\mathcal{S}_X| = 2m + 1$, razón por la cual cada punto posee una probabilidad elemental constante $P(X = k) = \dfrac{1}{2m+1}$.
	
	Analizamos la variable transformada $Z = X^2$:
	\begin{itemize}
		\item Para el punto de origen $z = 0$:
		La ecuación $X^2 = 0$ admite como única solución en el soporte a $X = 0$. Por lo tanto:
		\begin{align*}
			p_Z(0) = P(Z = 0) = P(X = 0) = \frac{1}{2m+1}.
		\end{align*}
		\item Para cualquier entero positivo del tipo $z = k^2$, con $k \in \set{1, 2, \dots, m}$:
		La ecuación cuadrática $X^2 = k^2$ admite exactamente dos soluciones enteras distintas dentro del soporte simétrico de $X$: $X = k$ y $X = -k$. En virtud de la aditividad de eventos mutuamente excluyentes, obtenemos:
		\begin{align*}
			p_Z(k^2) = P(Z = k^2) = P(X = k) + P(X = -k) = \frac{1}{2m+1} + \frac{1}{2m+1} = \frac{2}{2m+1}.
		\end{align*}
		\item Para todo valor $z \notin \set{0, 1^2, 2^2, \dots, m^2}$:
		La ecuación carece de preimágenes en $\mathcal{S}_X$, por lo que $p_Z(z) = 0$.
	\end{itemize}
	En conclusión, la función de masa queda formulada exactamente como:
	\begin{align*}
		p_Z(z) = \begin{cases}
			\dfrac{1}{2m+1} & \text{si } z = 0, \\[8pt]
			\dfrac{2}{2m+1} & \text{si } z \in \set{1^2, 2^2, \dots, m^2}, \\[8pt]
			0 & \text{en otro caso,}
		\end{cases}
	\end{align*}
	comprobando fehacientemente la afirmación teórica del enunciado.
\end{solucion}

\subsubsection*{4. Nivel Desafiante (Expertos / Paradojas)}

\begin{solucion}
	\label{sol:3.1.9}
	Para el Problema \ref{prob:3.1.9}:
	\begin{enumerate}
		\item La normalización de la distribución en el soporte infinito $\mathbb{Z}^+ = \set{1, 2, 3, \dots}$ se obtiene al sumar la serie infinita completa:
		\begin{align*}
			\sum_{k=1}^{\infty} f(k) = \sum_{k=1}^{\infty} \frac{1}{\zeta(s)\, k^s} = \frac{1}{\zeta(s)} \sum_{k=1}^{\infty} \frac{1}{k^s}.
		\end{align*}
		Dado que por definición analítica la función zeta de Riemann para $s > 1$ es convergente y satisface $\zeta(s) = \sum_{k=1}^{\infty} k^{-s}$, la serie se reduce a:
		\begin{align*}
			\sum_{k=1}^{\infty} f(k) = \frac{1}{\zeta(s)} \cdot \zeta(s) = 1.
		\end{align*}
		Esto certifica que la distribución de Zipf es una función de probabilidad estrictamente normalizada en su soporte infinito.
		
		\item Para acotar analíticamente la probabilidad de cola $P(X > M) = \sum_{k=M+1}^{\infty} \dfrac{1}{\zeta(s)\, k^s}$, utilizamos el criterio de comparación de sumas finitas e infinitas con integrales monótonas decrecientes. Como la función real $h(u) = u^{-s}$ es estrictamente decreciente para $u > 0$ y $s > 1$, se satisfacen las desigualdades fundamentales:
		\begin{align*}
			\int_{M+1}^{\infty} \frac{1}{u^s}\, du \le \sum_{k=M+1}^{\infty} \frac{1}{k^s} \le \int_{M}^{\infty} \frac{1}{u^s}\, du.
		\end{align*}
		Evaluamos la integral impropia genérica de la potencia negativa:
		\begin{align*}
			\int_{a}^{\infty} u^{-s}\, du = \left[ \frac{u^{-s+1}}{-s+1} \right]_{a}^{\infty} = \frac{a^{-(s-1)}}{s-1} = \frac{1}{(s-1)\, a^{s-1}}, \quad \text{para } a > 0.
		\end{align*}
		Sustituyendo los límites de integración correspondientes ($a = M+1$ por la izquierda y $a = M$ por la derecha), e incorporando el factor normalizador constante $1 / \zeta(s)$, se llega a la cota analítica exacta:
		\begin{align*}
			\frac{1}{\zeta(s)\,(s-1)\,(M+1)^{s-1}} \le P(X > M) \le \frac{1}{\zeta(s)\,(s-1)\,M^{s-1}}.
		\end{align*}
		Este resultado demuestra que las colas de la distribución de Pareto discreta decrecen según una ley de potencia polinomial ($\sim M^{-(s-1)}$), lo cual contrasta de forma radical con el decaimiento exponencial del tipo Poisson o geométrico.
	\end{enumerate}
\end{solucion}

\begin{solucion}
	\label{sol:3.1.10}
	Para el Problema \ref{prob:3.1.10}:
	\begin{enumerate}
		\item Para determinar la función de masa de probabilidad del tiempo hasta la primera falla combinada $T = \min(N_A, N_B)$, resulta matemáticamente más expedito trabajar con las probabilidades de supervivencia o colas superiores.
		La probabilidad de que un subsistema individual sobreviva más de $n-1$ ciclos discretos sin experimentar falla alguna es:
		\begin{align*}
			P(N_A > n-1) = (1 - p_A)^{n-1}, \quad P(N_B > n-1) = (1 - p_B)^{n-1}.
		\end{align*}
		Como los dos subsistemas operan de forma mutuamente independiente, el sistema completo sobrevive más allá del ciclo $n-1$ si y sólo si ambos subsistemas continúan operativos hasta ese momento:
		\begin{align*}
			P(T > n-1) &= P(\min(N_A, N_B) > n-1) = P(N_A > n-1 \text{ y } N_B > n-1) \\
			&= P(N_A > n-1) \cdot P(N_B > n-1) = \left[ (1 - p_A)(1 - p_B) \right]^{n-1}.
		\end{align*}
		Por la misma lógica, la probabilidad de supervivencia por encima de $n$ ciclos completos es:
		\begin{align*}
			P(T > n) = \left[ (1 - p_A)(1 - p_B) \right]^n.
		\end{align*}
		La función de masa de probabilidad en el ciclo exacto $n$ corresponde a la diferencia entre ambos sucesos:
		\begin{align*}
			f_T(n) = P(T = n) &= P(T > n-1) - P(T > n) \\
			&= \left[ (1 - p_A)(1 - p_B) \right]^{n-1} - \left[ (1 - p_A)(1 - p_B) \right]^n \\
			&= \left[ (1 - p_A)(1 - p_B) \right]^{n-1} \left( 1 - (1 - p_A)(1 - p_B) \right).
		\end{align*}
		
		\item Al expandir el término entre paréntesis:
		\begin{align*}
			1 - (1 - p_A)(1 - p_B) = 1 - (1 - p_A - p_B + p_A p_B) = p_A + p_B - p_A p_B.
		\end{align*}
		Definimos formalmente el parámetro de falla combinada del sistema en cada ciclo como:
		\begin{align*}
			p_{\text{eq}} = p_A + p_B - p_A p_B = 1 - (1 - p_A)(1 - p_B).
		\end{align*}
		Al sustituir $p_{\text{eq}}$ y su complemento $1 - p_{\text{eq}} = (1 - p_A)(1 - p_B)$ en la expresión deducida en el inciso anterior, obtenemos:
		\begin{align*}
			f_T(n) = (1 - p_{\text{eq}})^{n-1}\, p_{\text{eq}}, \quad n \in \set{1, 2, 3, \dots}.
		\end{align*}
		Queda demostrado que la variable aleatoria $T$ sigue de manera exacta una **distribución geométrica**, cuyo parámetro de probabilidad de falla es $p_{\text{eq}} = p_A + p_B - p_A p_B$.
	\end{enumerate}
\end{solucion}


% ==============================================================================
% SECCIÓN 03.02: FUNCIONES DE DISTRIBUCIÓN ACUMULADA (CDF DISCRETA)
% ==============================================================================
\subsection{Problemas sobre Funciones de Distribución Acumulada (CDF)}

\subsubsection*{1. Nivel Fundamental (Sencillos / Directos)}

\begin{problema}
	\label{prob:3.2.1}
	Considere un dado justo de seis caras equilibradas cuyo resultado se denota por la variable aleatoria discreta $X \in \set{1, 2, 3, 4, 5, 6}$.
	\begin{enumerate}
		\item Construya la expresión analítica exhaustiva de la función de distribución acumulada $F(x) = P(X \le x)$ para todo $x \in \mathbb{R}$.
		\item Calcule los valores numéricos de $F(2.8)$, $F(-1)$, $F(4)$ y $F(10)$.
	\end{enumerate}
\end{problema}

\begin{problema}
	\label{prob:3.2.2}
	Una variable aleatoria discreta $Y$ posee la siguiente función de distribución acumulada escalonada:
	\begin{align*}
		F(y) = \begin{cases}
			0 & \text{si } y < 0, \\
			0.15 & \text{si } 0 \le y < 2, \\
			0.55 & \text{si } 2 \le y < 5, \\
			0.85 & \text{si } 5 \le y < 8, \\
			1.00 & \text{si } y \ge 8.
		\end{cases}
	\end{align*}
	\begin{enumerate}
		\item Determine las probabilidades de intervalo $P(0 < Y \le 5)$ y $P(2 \le Y \le 8)$.
		\item Encuentre el soporte $\mathcal{S}_Y$ y construya la tabla completa de la función de masa de probabilidad $p_Y(y)$.
	\end{enumerate}
\end{problema}

\begin{problema}
	\label{prob:3.2.3}
	Sea $X$ una variable aleatoria discreta con función de distribución acumulada $F_X(x)$.
	\begin{enumerate}
		\item Explique rigurosamente por qué $F_X(x)$ debe ser una función no decreciente en toda la recta real ($\forall a \le b \implies F_X(a) \le F_X(b)$).
		\item Demuestre formalmente el comportamiento de continuidad por la derecha en cada punto del soporte ($x_k \in \mathcal{S}_X$), verificando que $\lim_{u \to x_k^+} F_X(u) = F_X(x_k)$.
	\end{enumerate}
\end{problema}


\subsubsection*{2. Nivel Operativo (Intermedios / Aplicados)}

\begin{problema}
	\label{prob:3.2.4}
	En un proceso automático de manufactura de circuitos integrados, el número de soldaduras defectuosas por oblea, $X$, presenta la función de masa de probabilidad:
	\begin{align*}
		p_X(x) = \begin{cases}
			0.40 & \text{si } x = 0, \\
			0.30 & \text{si } x = 1, \\
			0.20 & \text{si } x = 2, \\
			0.10 & \text{si } x = 3, \\
			0 & \text{en otro caso.}
		\end{cases}
	\end{align*}
	\begin{enumerate}
		\item Escriba la función de distribución acumulada $F_X(x)$ en forma de partes (función a trozos) para todo $x \in \mathbb{R}$.
		\item Si una oblea es rechazada cuando presenta $2$ o más defectos, calcule la probabilidad de rechazo utilizando exclusivamente los valores de $F_X(x)$.
	\end{enumerate}
\end{problema}

\begin{problema}
	\label{prob:3.2.5}
	Un equipo de ciberseguridad realiza intentos de penetración sobre un servidor web hasta encontrar la primera vulnerabilidad. Suponga que los intentos son independientes y que la probabilidad de éxito en cada intento individual es $p \in (0, 1)$. Sea $N$ el número total de intentos hasta el primer éxito inclusive.
	\begin{enumerate}
		\item Deduzca la expresión exacta en forma cerrada para la función de distribución acumulada $F_N(n) = P(N \le n)$ para cualquier entero $n \ge 1$.
		\item Si la probabilidad de vulnerar el servidor en un intento es $p = 0.20$, calcule el número mínimo de ensayos necesarios $n^*$ para que la probabilidad de éxito acumulado sea de al menos el $95\%$.
	\end{enumerate}
\end{problema}

\begin{problema}
	\label{prob:3.2.6}
	El cuantil de orden $\alpha \in (0, 1)$ para una variable aleatoria discreta $X$ se define convencionalmente como el ínfimo del conjunto de números reales donde la CDF supera o iguala a $\alpha$:
	\begin{align*}
		q_\alpha = \inf \set{x \in \mathbb{R} : F_X(x) \ge \alpha}.
	\end{align*}
	Considere que $X$ representa la demanda diaria de servidores en un centro de datos con distribución:
	$P(X=10)=0.10$, $P(X=20)=0.25$, $P(X=30)=0.35$, $P(X=40)=0.20$ y $P(X=50)=0.10$.
	\begin{enumerate}
		\item Calcule e interprete la mediana de la demanda diaria ($q_{0.50}$).
		\item Determine el cuantil de orden $0.85$ ($q_{0.85}$) y verifique su significado operativo en la gestión de capacidad del centro de datos.
	\end{enumerate}
\end{problema}


\subsubsection*{3. Nivel Analítico (Avanzados / Demostraciones)}

\begin{problema}
	\label{prob:3.2.7}
	Sea $X$ una variable aleatoria discreta cuyo soporte es un subconjunto de los números enteros ($\mathcal{S}_X \subset \mathbb{Z}$).
	\begin{enumerate}
		\item Demuestre rigurosamente que para cualesquiera dos enteros $a, b \in \mathbb{Z}$ con $a \le b$, la probabilidad del intervalo cerrado se expresa exactamente como:
		\begin{align*}
			P(a \le X \le b) = F_X(b) - F_X(a - 1).
		\end{align*}
		\item Demuestre que si $X$ es entera, la magnitud del salto en cualquier entero $k \in \mathbb{Z}$ satisface la relación de diferencia finita $p_X(k) = \Delta F_X(k) = F_X(k) - F_X(k-1)$.
	\end{enumerate}
\end{problema}

\begin{problema}
	\label{prob:3.2.8}
	Considere una variable aleatoria discreta $X$ con soporte $\mathcal{S}_X = \set{x_1, x_2, \dots}$ ordenado de manera estrictamente creciente ($x_1 < x_2 < \dots$) y función de distribución $F_X(x)$. Sea $g: \mathbb{R} \to \mathbb{R}$ una función estrictamente monótona creciente en el soporte de $X$, y definamos la nueva variable $Y = g(X)$.
	\begin{enumerate}
		\item Demuestre teóricamente que la función de distribución de la variable transformada, $F_Y(y)$, preserva la estructura escalonada de $F_X$ y se relaciona mediante:
		\begin{align*}
			F_Y(y) = F_X\left( g^{-1}(y) \right) \quad \text{para todo } y \in g(\mathcal{S}_X).
		\end{align*}
		\item Si en cambio $h: \mathbb{R} \to \mathbb{R}$ es una función estrictamente decreciente y $Z = h(X)$, deduzca la fórmula analítica para $F_Z(z)$ en términos de $F_X$.
	\end{enumerate}
\end{problema}


\subsubsection*{4. Nivel Desafiante (Expertos / Paradojas)}

\begin{problema}
	\label{prob:3.2.9}
	En la teoría de la medida y probabilidad avanzada, se distinguen las distribuciones discretas puras de las distribuciones absolutamente continuas y de las singulares continuas (como la escalera de Cantor).
	\begin{enumerate}
		\item Demuestre analíticamente que la función de distribución acumulada $F(x)$ de cualquier variable aleatoria discreta pura se puede expresar como una serie convergente de funciones escalón unitarias de Heaviside $H(x - x_k)$:
		\begin{align*}
			F(x) = \sum_{k} p_k\, H(x - x_k), \quad \text{donde } H(t) = \begin{cases} 1 & \text{si } t \ge 0, \\ 0 & \text{si } t < 0. \end{cases}
		\end{align*}
		\item Explique por qué la derivada ordinaria $F'(x)$ es idéntica a cero en casi todas partes ($\text{c.t.p.}$) respecto a la medida de Lebesgue en $\mathbb{R}$, y cómo la teoría de distribuciones (o funciones generalizadas) recupera la masa puntual en el sentido de la delta de Dirac: $f(x) = \sum p_k \delta(x - x_k)$.
	\end{enumerate}
\end{problema}

\begin{problema}
	\label{prob:3.2.10}
	Considere un par de variables aleatorias discretas $(X, Y)$ que representan el tráfico de paquetes de entrada y salida en un enrutador óptico en intervalos discretos de tiempo, con función de distribución acumulada conjunta discreta $F_{X,Y}(x, y) = P(X \le x, Y \le y)$.
	\begin{enumerate}
		\item Demuestre algebraicamente que la probabilidad de que el par $(X, Y)$ caiga dentro de un rectángulo semiabierto $(x_1, x_2] \times (y_1, y_2]$ se expresa mediante las segundas diferencias de la CDF conjunta:
		\begin{align*}
			P(x_1 < X \le x_2, y_1 < Y \le y_2) = F_{X,Y}(x_2, y_2) - F_{X,Y}(x_1, y_2) - F_{X,Y}(x_2, y_1) + F_{X,Y}(x_1, y_1).
		\end{align*}
		\item Demuestre que para cualquier par de variables discretas, la CDF conjunta está fundamentalmente acotada en todo punto $(x, y) \in \mathbb{R}^2$ por las cotas de Fréchet-Hoeffding:
		\begin{align*}
			\max\left( 0, F_X(x) + F_Y(y) - 1 \right) \le F_{X,Y}(x, y) \le \min\left( F_X(x), F_Y(y) \right).
		\end{align*}
	\end{enumerate}
\end{problema}


% ==============================================================================
% SUGERENCIAS PARA LA SECCIÓN 03.02
% ==============================================================================
\subsection{Sugerencias para los problemas de la Sección 03.02}

\begin{sugerencia}
	\label{sug:3.2.1}
	Para el Problema \ref{prob:3.2.1}: Divida la recta real en 7 intervalos disjuntos: $(-\infty, 1)$, $[1, 2)$, $[2, 3)$, $[3, 4)$, $[4, 5)$, $[5, 6)$ y $[6, \infty)$. En cada intervalo, sume las probabilidades $f(u) = 1/6$ de las caras que sean menores o iguales a $x$.
\end{sugerencia}

\begin{sugerencia}
	\label{sug:3.2.2}
	Para el Problema \ref{prob:3.2.2}: Recuerde la identidad fundamental $P(a < Y \le b) = F(b) - F(a)$. Para reconstruir la PMF, identifique los puntos donde $F(y)$ experimenta saltos e inspeccione las diferencias $F(y_k) - F(y_k^-)$.
\end{sugerencia}

\begin{sugerencia}
	\label{sug:3.2.3}
	Para el Problema \ref{prob:3.2.3}: Utilice el hecho de que si $a \le b$, el evento $\set{\omega : X(\omega) \le a}$ es un subconjunto del evento $\set{\omega : X(\omega) \le b}$. Aplique la monotonía de la medida de probabilidad y el axioma de continuidad por arriba sobre secuencias decrecientes de intervalos.
\end{sugerencia}

\begin{sugerencia}
	\label{sug:3.2.4}
	Para el Problema \ref{prob:3.2.4}: Sume progresivamente las probabilidades $p_X(x)$ para $x=0, 1, 2, 3$. Exprese la probabilidad de rechazo como la probabilidad del complemento de aceptar, es decir, $P(X \ge 2) = 1 - P(X \le 1) = 1 - F_X(1)$.
\end{sugerencia}

\begin{sugerencia}
	\label{sug:3.2.5}
	Para el Problema \ref{prob:3.2.5}: El evento complementario $\set{N > n}$ equivale a que los primeros $n$ ensayos consecutivos hayan resultado en fracasos independientes, cuya probabilidad es $(1-p)^n$. Luego, plantee la inecuación $1 - (1-p)^n \ge 0.95$ y resuelva usando logaritmos naturales.
\end{sugerencia}

\begin{sugerencia}
	\label{sug:3.2.6}
	Para el Problema \ref{prob:3.2.6}: Calcule la tabla de acumulaciones $F_X(x)$ para $x \in \set{10, 20, 30, 40, 50}$. Inspeccione cuál es el menor valor $x_k$ en el cual la probabilidad acumulada alcanza o sobrepasa $0.50$ (para la mediana) y $0.85$ (para el percentil superior).
\end{sugerencia}

\begin{sugerencia}
	\label{sug:3.2.7}
	Para el Problema \ref{prob:3.2.7}: Si $X$ solo toma valores enteros, el conjunto de números enteros en el intervalo cerrado $[a, b]$ es exactamente la diferencia de conjuntos entre $(-\infty, b] \cap \mathbb{Z}$ y $(-\infty, a-1] \cap \mathbb{Z}$. Exprese esto mediante probabilidades algebraicas.
\end{sugerencia}

\begin{sugerencia}
	\label{sug:3.2.8}
	Para el Problema \ref{prob:3.2.8}: Si $g$ es estrictamente creciente en la recta real o sobre el soporte de $X$, la desigualdad $Y \le y \iff g(X) \le y$ es lógicamente equivalente a $X \le g^{-1}(y)$. Para el caso decreciente, invierta el sentido de la desigualdad y preste atención estricta al límite por la izquierda.
\end{sugerencia}

\begin{sugerencia}
	\label{sug:3.2.9}
	Para el Problema \ref{prob:3.2.9}: Una función escalón de Heaviside $H(x - x_k)$ vale $0$ para $x < x_k$ y $1$ para $x \ge x_k$. Verifique que al ponderar cada escalón por $p_k = P(X = x_k)$ se reproduce exactamente la suma de masa $\sum_{x_k \le x} p_k$. Para la derivada, recuerde que una constante a trozos tiene pendiente cero en todo el interior de cada subintervalo continuo.
\end{sugerencia}

\begin{sugerencia}
	\label{sug:3.2.10}
	Para el Problema \ref{prob:3.2.10}: Para el inciso 1, dibuje el rectángulo en el plano cartesiano $(x, y)$ y exprese la masa interna restando las CDFs de las regiones exteriores inframarginales. Para el inciso 2, utilice el hecho de que la probabilidad conjunta nunca puede exceder ninguna de las dos marginales ni ser inferior a la suma combinada menos la unidad.
\end{sugerencia}


% ==============================================================================
% SOLUCIONES PARA LA SECCIÓN 03.02
% ==============================================================================
\subsection{Soluciones a los problemas de la Sección 03.02}

\begin{solucion}
	\label{sol:3.2.1}
	Para el Problema \ref{prob:3.2.1}:
	\begin{enumerate}
		\item Para un dado justo con soporte $\mathcal{S}_X = \set{1, 2, 3, 4, 5, 6}$, cada cara tiene probabilidad puntual constante $f(k) = \frac{1}{6}$.
		La función de distribución acumulada se define analíticamente sobre todo el eje real como la suma de las masas a la izquierda o en la frontera de $x$:
		\begin{align*}
			F(x) = P(X \le x) = \sum_{k \le x} f(k) = \begin{cases}
				0 & \text{si } x < 1, \\[3pt]
				\frac{1}{6} & \text{si } 1 \le x < 2, \\[3pt]
				\frac{2}{6} = \frac{1}{3} & \text{si } 2 \le x < 3, \\[3pt]
				\frac{3}{6} = \frac{1}{2} & \text{si } 3 \le x < 4, \\[3pt]
				\frac{4}{6} = \frac{2}{3} & \text{si } 4 \le x < 5, \\[3pt]
				\frac{5}{6} & \text{si } 5 \le x < 6, \\[3pt]
				1 & \text{si } x \ge 6.
			\end{cases}
		\end{align*}
		Esta función escalonada presenta saltos de magnitud idéntica a $\frac{1}{6}$ exactamente en cada entero del soporte.
		
		\item Evaluando numéricamente $F(x)$ en cada punto requerido observando el intervalo correspondiente de la función a trozos:
		\begin{itemize}
			\item Para $x = 2.8$: Como $2 \le 2.8 < 3$, el valor acumulado es $F(2.8) = \dfrac{2}{6} = \dfrac{1}{3} \approx 0.3333$.
			\item Para $x = -1$: Como $-1 < 1$, se encuentra en la región de acumulación nula, por lo que $F(-1) = 0$.
			\item Para $x = 4$: Al ser exactamente un punto de discontinuidad (salto), aplicamos la continuidad por la derecha ($4 \le 4 < 5$), obteniendo $F(4) = \dfrac{4}{6} = \dfrac{2}{3} \approx 0.6667$.
			\item Para $x = 10$: Como $10 \ge 6$, el acumulado ha saturado el espacio muestral total, por lo que $F(10) = 1$.
		\end{itemize}
	\end{enumerate}
\end{solucion}

\begin{solucion}
	\label{sol:3.2.2}
	Para el Problema \ref{prob:3.2.2}:
	\begin{enumerate}
		\item Utilizando la relación analítica para intervalos semiabiertos en variables aleatorias cualesquiera:
		\begin{align*}
			P(a < Y \le b) = P(Y \le b) - P(Y \le a) = F(b) - F(a).
		\end{align*}
		Evaluando cada intervalo solicitado directamente sobre los valores explícitos de la CDF:
		\begin{align*}
			P(0 < Y \le 5) &= F(5) - F(0) = 0.85 - 0.15 = 0.70. \\
			P(2 \le Y \le 8) &= P(Y = 2) + P(2 < Y \le 8) \\
			&= \left[ F(2) - F(2^-) \right] + \left[ F(8) - F(2) \right] = F(8) - F(2^-).
		\end{align*}
		Como el límite por la izquierda en el salto es $F(2^-) = \lim_{u \to 2^-} F(u) = 0.15$, obtenemos:
		\begin{align*}
			P(2 \le Y \le 8) = 1.00 - 0.15 = 0.85.
		\end{align*}
		
		\item El soporte $\mathcal{S}_Y$ está conformado por los puntos exactos de discontinuidad (saltos) de $F(y)$, ya que son los únicos números reales donde se concentra masa positiva ($f(y_k) > 0$).
		Observando los cambios de meseta en la función a trozos, identificamos que los saltos ocurren en $y \in \set{0, 2, 5, 8}$. Por tanto:
		\begin{align*}
			\mathcal{S}_Y = \set{0, 2, 5, 8}.
		\end{align*}
		Las magnitudes de las masas puntuales se obtienen por diferencia finita $p_Y(y_k) = F(y_k) - F(y_k^-)$:
		\begin{align*}
			p_Y(0) &= F(0) - F(0^-) = 0.15 - 0 = 0.15, \\
			p_Y(2) &= F(2) - F(2^-) = 0.55 - 0.15 = 0.40, \\
			p_Y(5) &= F(5) - F(5^-) = 0.85 - 0.55 = 0.30, \\
			p_Y(8) &= F(8) - F(8^-) = 1.00 - 0.85 = 0.15.
		\end{align*}
		La tabla completa de la función de masa de probabilidad verifica el axioma de normalización unitaria:
		\begin{center}
			\small
			\begin{tabular}{c|cccc|c}
				$y$ & $0$ & $2$ & $5$ & $8$ & Suma \\
				\hline
				$p_Y(y)$ & $0.15$ & $0.40$ & $0.30$ & $0.15$ & $1.00$
			\end{tabular}
		\end{center}
	\end{enumerate}
\end{solucion}

\begin{solucion}
	\label{sol:3.2.3}
	Para el Problema \ref{prob:3.2.3}:
	\begin{enumerate}
		\item Sean $a, b \in \mathbb{R}$ tales que $a \le b$. Por definición formal de relación de orden sobre los números reales, cualquier realización elemental $\omega \in \Omega$ que satisfaga la condición de pertenencia al evento anterior ($X(\omega) \le a$) satisface automáticamente la condición de pertenencia al evento posterior ($X(\omega) \le b$).
		En términos de teoría de conjuntos medibles:
		\begin{align*}
			A = \set{\omega \in \Omega : X(\omega) \le a} \subseteq B = \set{\omega \in \Omega : X(\omega) \le b}.
		\end{align*}
		Por la propiedad de monotonía de cualquier medida de probabilidad legítima (Axiomas de Kolmogorov), si $A \subseteq B$, entonces sus medidas cumplen $P(A) \le P(B)$. En consecuencia directa:
		\begin{align*}
			F_X(a) = P(X \le a) \le P(X \le b) = F_X(b).
		\end{align*}
		Por lo tanto, la función de distribución acumulada es estrictamente monotónica no decreciente en todo su dominio.
		
		\item Para demostrar la continuidad por la derecha en un punto arbitrario $x \in \mathbb{R}$, considere una sucesión decreciente de números reales $\set{x_n}_{n=1}^\infty$ que converja hacia $x$ desde la derecha ($x_n \downarrow x$, es decir, $x_1 > x_2 > \dots > x_n > \dots$ con $\lim_{n \to \infty} x_n = x$).
		Definamos la familia de eventos medibles decrecientes:
		\begin{align*}
			E_n = \set{\omega \in \Omega : X(\omega) \le x_n}.
		\end{align*}
		Como $x_1 > x_2 > \dots$, se tiene la inclusión anidada $E_1 \supseteq E_2 \supseteq \dots \supseteq E_n \supseteq \dots$. La intersección numerable de todos estos conjuntos es precisamente el evento en el límite de convergencia:
		\begin{align*}
			\bigcap_{n=1}^\infty E_n = \set{\omega \in \Omega : X(\omega) \le x}.
		\end{align*}
		Por la propiedad fundamental de continuidad de la medida de probabilidad por arriba sobre secuencias anidadas de eventos medibles:
		\begin{align*}
			\lim_{n \to \infty} F_X(x_n) = \lim_{n \to \infty} P(E_n) = P\left( \bigcap_{n=1}^\infty E_n \right) = P(X \le x) = F_X(x).
		\end{align*}
		Como este resultado es independiente de la sucesión decreciente particular $\set{x_n}$ elegida, se demuestra algebraicamente que $\lim_{u \to x^+} F_X(u) = F_X(x)$ para todo $x \in \mathbb{R}$.
	\end{enumerate}
\end{solucion}

\begin{solucion}
	\label{sol:3.2.4}
	Para el Problema \ref{prob:3.2.4}:
	\begin{enumerate}
		\item Acumulando progresivamente las probabilidades puntuales dadas por $p_X(x)$ sobre los intervalos definidos por el soporte $\mathcal{S}_X = \set{0, 1, 2, 3}$:
		\begin{align*}
			F_X(x) = \sum_{u \le x} p_X(u) = \begin{cases}
				0 & \text{si } x < 0, \\[3pt]
				0.40 & \text{si } 0 \le x < 1, \\[3pt]
				0.40 + 0.30 = 0.70 & \text{si } 1 \le x < 2, \\[3pt]
				0.70 + 0.20 = 0.90 & \text{si } 2 \le x < 3, \\[3pt]
				0.90 + 0.10 = 1.00 & \text{si } x \ge 3.
			\end{cases}
		\end{align*}
		
		\item La probabilidad operativa de rechazo en manufactura corresponde al suceso de que la oblea presente dos o más soldaduras defectuosas ($X \ge 2$).
		Utilizando el principio del evento complementario y evaluando en la CDF:
		\begin{align*}
			P(\text{Rechazo}) &= P(X \ge 2) = 1 - P(X < 2).
		\end{align*}
		Dado que la variable solo toma valores enteros discretos, el suceso estricto $X < 2$ es idéntico al suceso acotado $X \le 1$. Por ende:
		\begin{align*}
			P(\text{Rechazo}) = 1 - F_X(1) = 1 - 0.70 = 0.30.
		\end{align*}
		Existe un $30\%$ de probabilidad exacta de que una oblea sea descartada por fallas de soldadura en cada lote de control de calidad.
	\end{enumerate}
\end{solucion}

\begin{solucion}
	\label{sol:3.2.5}
	Para el Problema \ref{prob:3.2.5}:
	\begin{enumerate}
		\item Sea $N$ el número de intentos independientes hasta el primer éxito, el cual sigue por construcción una distribución geométrica sobre $\mathcal{S}_N = \set{1, 2, 3, \dots}$ con probabilidades puntuales $P(N = k) = (1-p)^{k-1}\, p$.
		Para deducir la forma cerrada elegante de la función de distribución acumulada $F_N(n) = P(N \le n)$ para $n \ge 1$, recurrimos al evento complementario: el suceso $\set{N > n}$ significa formal y físicamente que los primeros $n$ intentos consecutivos resultaron infructuosos (fracasos).
		Como los intentos son eventos independientes entre sí y la probabilidad de fallo por intento es $1 - p$:
		\begin{align*}
			P(N > n) = P(\text{Fracaso en ensayo } 1 \text{ y } \dots \text{ y Fracaso en ensayo } n) = (1 - p)^n.
		\end{align*}
		Tomando el complemento para encontrar la probabilidad de acumulación o éxito en los primeros $n$ intentos:
		\begin{align*}
			F_N(n) = P(N \le n) = 1 - P(N > n) = 1 - (1 - p)^n, \quad \text{para todo } n \in \set{1, 2, 3, \dots}.
		\end{align*}
		
		\item Evaluando para un parámetro de penetración de ciberseguridad $p = 0.20$, el complemento de fallo es $1 - p = 0.80$. Requerimos encontrar el entero mínimo $n^*$ tal que:
		\begin{align*}
			F_N(n) \ge 0.95 \implies 1 - (0.80)^n \ge 0.95.
		\end{align*}
		Despejando el término exponencial:
		\begin{align*}
			(0.80)^n \le 0.05.
		\end{align*}
		Aplicando el logaritmo natural a ambos miembros de la inecuación (recordando que $\ln(0.80) < 0$, por lo que el sentido de la desigualdad se invierte al dividir):
		\begin{align*}
			n \ln(0.80) \le \ln(0.05) \implies n \ge \frac{\ln(0.05)}{\ln(0.80)}.
		\end{align*}
		Calculando numéricamente con precisión:
		\begin{align*}
			n \ge \frac{-2.99573}{-0.22314} \approx 13.425.
		\end{align*}
		Como el número de ensayos es estrictamente un número entero positivo, aplicamos la función techo:
		\begin{align*}
			n^* = \lceil 13.425 \rceil = 14.
		\end{align*}
		El equipo de ciberseguridad debe ejecutar al menos $14$ intentos independientes para garantizar con una confianza acumulada del $95\%$ o superior la detección de la primera vulnerabilidad.
	\end{enumerate}
\end{solucion}

\begin{solucion}
	\label{sol:3.2.6}
	Para el Problema \ref{prob:3.2.6}:
	\begin{enumerate}
		\item Primero, construimos la tabla completa de la función de distribución acumulada de la demanda diaria en el centro de datos sumando progresivamente las probabilidades:
		\begin{center}
			\small
			\begin{tabular}{c|ccccc}
				$x$ (Demanda en Servidores) & $10$ & $20$ & $30$ & $40$ & $50$ \\
				\hline
				$p_X(x)$ & $0.10$ & $0.25$ & $0.35$ & $0.20$ & $0.10$ \\
				$F_X(x)$ & $0.10$ & $0.35$ & $0.70$ & $0.90$ & $1.00$
			\end{tabular}
		\end{center}
		La mediana o cuantil $50$ ($q_{0.50}$) se obtiene identificando el menor valor de demanda para el cual el acumulado $F_X(x)$ alcanza o supera a $0.50$:
		\begin{align*}
			q_{0.50} = \inf \set{x \in \mathcal{S}_X : F_X(x) \ge 0.50}.
		\end{align*}
		Inspeccionando la tabla: $F_X(10) = 0.10 < 0.50$, $F_X(20) = 0.35 < 0.50$, y $F_X(30) = 0.70 \ge 0.50$.
		Por lo tanto:
		\begin{align*}
			q_{0.50} = 30 \text{ servidores}.
		\end{align*}
		Interpretación: El $50\%$ o más de los días hábiles operativos del centro de datos requieren como máximo 30 servidores activos para abastecer toda la demanda de cómputo diario, y el otro $50\%$ de los días requieren como mínimo dicha cantidad.
		
		\item Para calcular el cuantil de orden $0.85$ ($q_{0.85}$):
		\begin{align*}
			q_{0.85} = \inf \set{x \in \mathcal{S}_X : F_X(x) \ge 0.85}.
		\end{align*}
		Observando las acumulaciones sucesivas: $F_X(30) = 0.70 < 0.85$, mientras que $F_X(40) = 0.90 \ge 0.85$.
		En consecuencia directa:
		\begin{align*}
			q_{0.85} = 40 \text{ servidores}.
		\end{align*}
		Significado operativo: Si el equipo de ingeniería de infraestructura asigna y aprovisiona una capacidad estática constante de 40 servidores por día, el centro de datos operará sin saturación de recursos ni pérdidas de servicio en el $90\%$ de los días (satisfaciendo con amplio margen el requerimiento operativo de un nivel de servicio del $85\%$).
	\end{enumerate}
\end{solucion}

\begin{solucion}
	\label{sol:3.2.7}
	Para el Problema \ref{prob:3.2.7}:
	\begin{enumerate}
		\item Sean $a, b \in \mathbb{Z}$ tales que $a \le b$. Consideremos la descomposición en conjuntos algebraicos disjuntos del suceso acotado superiormente por $b$:
		\begin{align*}
			\set{\omega \in \Omega : X(\omega) \le b} = \set{\omega \in \Omega : X(\omega) \le a - 1} \cup \set{\omega \in \Omega : a \le X(\omega) \le b}.
		\end{align*}
		Esta igualdad de conjuntos es algebraicamente exacta bajo la hipótesis de que $X$ es una variable con soporte exclusivamente entero ($\mathcal{S}_X \subset \mathbb{Z}$), dado que no existen números enteros estrictamente comprendidos en el intervalo abierto $(a-1, a)$.
		Como la unión anterior es sobre eventos mutuamente excluyentes, el tercer axioma de Kolmogorov establece la aditividad de la probabilidad:
		\begin{align*}
			P(X \le b) = P(X \le a - 1) + P(a \le X \le b).
		\end{align*}
		Sustituido en términos de la definición de la CDF ($F_X(x) = P(X \le x)$) y despejando la probabilidad del intervalo cerrado:
		\begin{align*}
			P(a \le X \le b) = F_X(b) - F_X(a - 1).
		\end{align*}
		Queda rigurosamente demostrada la fórmula de diferencia finita para intervalos sobre el conjunto de los enteros.
		
		\item Para encontrar la magnitud del salto en un punto entero cualquiera $k \in \mathbb{Z}$, evaluamos la fórmula del intervalo cerrado obtenida en el inciso previo haciendo $a = b = k$:
		\begin{align*}
			p_X(k) = P(X = k) = P(k \le X \le k) = F_X(k) - F_X(k - 1) = \Delta F_X(k).
		\end{align*}
		Esta relación demuestra explícitamente que la función de masa de probabilidad en cualquier entero es idéntica al operador de primera diferencia finita hacia atrás aplicado a la función de distribución acumulada.
	\end{enumerate}
\end{solucion}

\begin{solucion}
	\label{sol:3.2.8}
	Para el Problema \ref{prob:3.2.8}:
	\begin{enumerate}
		\item Sea $g: \mathbb{R} \to \mathbb{R}$ una función estrictamente monótona creciente en el soporte de $X$. Por la definición formal de monotonía estricta creciente, para cualquier par de números reales $x, y$, la relación de desigualdad se preserva bajo la aplicación del operador inverso:
		\begin{align*}
			g(x) \le y \iff x \le g^{-1}(y).
		\end{align*}
		Definiendo la variable aleatoria transformada $Y = g(X)$, calculamos su función de distribución acumulada en cualquier punto del recorrido $y \in \mathbb{R}$:
		\begin{align*}
			F_Y(y) = P(Y \le y) = P(g(X) \le y) = P(X \le g^{-1}(y)) = F_X\left( g^{-1}(y) \right).
		\end{align*}
		Como la composición de una función escalonada por la derecha ($F_X$) con una función biyectiva creciente ($g^{-1}$) produce una nueva función escalonada con puntos de salto exactamente en $y_k = g(x_k)$, se comprueba formalmente la conservación de la estructura discreta y la identidad teórica.
		
		\item Sea $h: \mathbb{R} \to \mathbb{R}$ una función estrictamente monótona decreciente en el soporte de $X$, y definamos $Z = h(X)$.
		Al aplicar una transformación estrictamente decreciente, el sentido de la desigualdad matemática se invierte en el dominio de la preimagen:
		\begin{align*}
			h(X) \le z \iff X \ge h^{-1}(z).
		\end{align*}
		Evaluando la probabilidad acumulada para la variable $Z$:
		\begin{align*}
			F_Z(z) = P(Z \le z) = P(h(X) \le z) = P(X \ge h^{-1}(z)).
		\end{align*}
		Para expresar el suceso $X \ge h^{-1}(z)$ a partir de la función de distribución acumulada (la cual por definición acumula hacia la izquierda estrictamente como $P(X \le x)$), utilizamos la probabilidad del evento complementario y su límite por la izquierda:
		\begin{align*}
			P(X \ge h^{-1}(z)) &= 1 - P(X < h^{-1}(z)) \\
			&= 1 - \lim_{u \to (h^{-1}(z))^-} F_X(u) = 1 - F_X\left( \left[h^{-1}(z)\right]^- \right).
		\end{align*}
		En conclusión, para una transformación decreciente, la función de distribución acumulada de la nueva variable discreta se rige algebraicamente por la fórmula exacto-analítica:
		\begin{align*}
			F_Z(z) = 1 - F_X\left( \left[h^{-1}(z)\right]^- \right).
		\end{align*}
	\end{enumerate}
\end{solucion}

\begin{solucion}
	\label{sol:3.2.9}
	Para el Problema \ref{prob:3.2.9}:
	\begin{enumerate}
		\item Sea $X$ una variable aleatoria discreta pura cuyo soporte medible es $\mathcal{S}_X = \set{x_k}_{k \in \mathcal{K}}$, donde el índice $\mathcal{K}$ es a lo sumo infinito numerable y cada punto concentra probabilidad $p_k = P(X = x_k) > 0$ con $\sum_k p_k = 1$.
		La función de distribución acumulada en cualquier punto $x \in \mathbb{R}$ se escribe como la suma de las masas de probabilidad sobre los puntos del soporte menores o iguales a $x$:
		\begin{align*}
			F(x) = \sum_{x_k \le x} p_k.
		\end{align*}
		Considere ahora la función escalón unitario de Heaviside $H(t)$, que satisface:
		\begin{align*}
			H(x - x_k) = \begin{cases}
				1 & \text{si } x - x_k \ge 0 \iff x \ge x_k, \\
				0 & \text{si } x - x_k < 0 \iff x < x_k.
			\end{cases}
		\end{align*}
		Multiplicando cada término de la serie infinita por el indicador $H(x - x_k)$, todo sumando donde $x < x_k$ se anula automáticamente (por valer cero el escalón), mientras que todo sumando donde $x \ge x_k$ contribuye íntegramente con su masa $p_k \cdot 1 = p_k$:
		\begin{align*}
			\sum_{k \in \mathcal{K}} p_k\, H(x - x_k) = \sum_{k : x_k \le x} p_k \cdot 1 + \sum_{k : x_k > x} p_k \cdot 0 = \sum_{x_k \le x} p_k = F(x).
		\end{align*}
		Dado que por el segundo axioma de Kolmogorov la serie de probabilidades es absolutamente convergente ($\sum |p_k| = \sum p_k = 1 < \infty$), esta serie de funciones de Heaviside converge uniformemente en todo $\mathbb{R}$, demostrando la representación estructural discreta.
		
		\item En el sentido del cálculo diferencial ordinario (análisis real sobre la medida de Lebesgue en $\mathbb{R}$), el complemento del soporte numerable $\mathbb{R} \setminus \mathcal{S}_X$ es un conjunto abierto de medida total, constituido por una unión numerable de intervalos disjuntos dentro de los cuales la función escalonada $F(x)$ es local y absolutamente constante (mesetas horizontales).
		Por consiguiente, la derivada ordinaria en cada meseta es idénticamente nula:
		\begin{align*}
			F'(x) = 0 \quad \text{para todo } x \notin \mathcal{S}_X.
		\end{align*}
		Como la cardinalidad de los puntos de discontinuidad $\mathcal{S}_X$ es numerable (medida de Lebesgue igual a cero, $\mu_L(\mathcal{S}_X) = 0$), se concluye en el análisis de variables reales que $F'(x) = 0$ casi en todas partes ($\text{c.t.p.}$).
		
		Sin embargo, para preservar la coherencia y rigurosidad analítica en ingeniería y física teórica sin perder la masa de probabilidad total, se invoca la **teoría de distribuciones (o funciones generalizadas de Sobolev-Schwartz)**. En este marco topológico, la derivada distribucional de la función escalón unitaria de Heaviside $H(t)$ se define como la medida de masa puntual de Dirac centrada en el origen, $\delta(t)$, cuya integral satisface $\int_{-\infty}^\infty \delta(t) dt = 1$.
		Al derivar término a término la serie convergente deducida en el inciso anterior en el sentido de las distribuciones, obtenemos:
		\begin{align*}
			f(x) = \frac{d}{dx} F(x) = \frac{d}{dx} \left[ \sum_{k \in \mathcal{K}} p_k\, H(x - x_k) \right] = \sum_{k \in \mathcal{K}} p_k\, \frac{d}{dx} H(x - x_k) = \sum_{k \in \mathcal{K}} p_k\, \delta(x - x_k).
		\end{align*}
		Esta hermosa identidad generalizada reconcilia el cálculo diferencial continuo con la probabilidad discreta, demostrando que la densidad generalizada de una variable aleatoria discreta pura es una superposición lineal de impulsos de Dirac ponderados por sus probabilidades de masa.
	\end{enumerate}
\end{solucion}

\begin{solucion}
	\label{sol:3.2.10}
	Para el Problema \ref{prob:3.2.10}:
	\begin{enumerate}
		\item Sea el rectángulo semiabierto $\mathcal{R} = (x_1, x_2] \times (y_1, y_2]$ en el plano cartesiano con $x_1 < x_2$ y $y_1 < y_2$. La probabilidad conjunta de que el vector aleatorio $(X, Y)$ caiga dentro de este recinto se denota por $P((X,Y) \in \mathcal{R})$.
		Por la geometría de los conjuntos medibles en $\mathbb{R}^2$, podemos expresar el cuadrante infinito inferior izquierdo con vértice en $(x_2, y_2)$ como la unión disjunta de cuatro regiones rectangulares:
		\begin{align*}
			(-\infty, x_2] \times (-\infty, y_2] &= \mathcal{R} \cup \left[ (-\infty, x_1] \times (y_1, y_2] \right] \\
			&\quad \cup \left[ (x_1, x_2] \times (-\infty, y_1] \right] \cup \left[ (-\infty, x_1] \times (-\infty, y_1] \right].
		\end{align*}
		Evaluando la medida de probabilidad conjunta en el cuadrante total exterior (que por definición formal constituye $F_{X,Y}(x_2, y_2)$):
		\begin{align*}
			F_{X,Y}(x_2, y_2) &= P((X,Y) \in \mathcal{R}) + P(X \le x_1, y_1 < Y \le y_2) \\
			&\quad + P(x_1 < X \le x_2, Y \le y_1) + P(X \le x_1, Y \le y_1).
		\end{align*}
		Descomponiendo los términos intermedios mediante restas de funciones de distribución:
		\begin{align*}
			P(X \le x_1, y_1 < Y \le y_2) &= P(X \le x_1, Y \le y_2) - P(X \le x_1, Y \le y_1) \\
			&= F_{X,Y}(x_1, y_2) - F_{X,Y}(x_1, y_1), \\[4pt]
			P(x_1 < X \le x_2, Y \le y_1) &= P(X \le x_2, Y \le y_1) - P(X \le x_1, Y \le y_1) \\
			&= F_{X,Y}(x_2, y_1) - F_{X,Y}(x_1, y_1).
		\end{align*}
		Sustituyendo ambas identidades algebraicas en la ecuación de conservación original y agrupando términos:
		\begin{align*}
			F_{X,Y}(x_2, y_2) &= P((X,Y) \in \mathcal{R}) + \left[ F_{X,Y}(x_1, y_2) - F_{X,Y}(x_1, y_1) \right] \\
			&\quad + \left[ F_{X,Y}(x_2, y_1) - F_{X,Y}(x_1, y_1) \right] + F_{X,Y}(x_1, y_1) \\
			&= P((X,Y) \in \mathcal{R}) + F_{X,Y}(x_1, y_2) + F_{X,Y}(x_2, y_1) - F_{X,Y}(x_1, y_1).
		\end{align*}
		Despejando finalmente la masa de probabilidad contenida en el interior del rectángulo $\mathcal{R}$:
		\begin{align*}
			P(x_1 < X \le x_2, y_1 < Y \le y_2) = F_{X,Y}(x_2, y_2) - F_{X,Y}(x_1, y_2) - F_{X,Y}(x_2, y_1) + F_{X,Y}(x_1, y_1).
		\end{align*}
		Esta expresión representa la segunda diferencia mixta o enésima variación del operador acumulativo bidimensional.
		
		\item Para demostrar la cota superior de Fréchet-Hoeffding en cualquier punto $(x, y) \in \mathbb{R}^2$, notemos que el evento bivariado de intersección está contenido trivialmente dentro de cada uno de los dos eventos marginales por separado:
		\begin{align*}
			\set{\omega \in \Omega : X(\omega) \le x \text{ y } Y(\omega) \le y} \subseteq \set{\omega \in \Omega : X(\omega) \le x}, \\
			\set{\omega \in \Omega : X(\omega) \le x \text{ y } Y(\omega) \le y} \subseteq \set{\omega \in \Omega : Y(\omega) \le y}.
		\end{align*}
		Por la monotonía de la medida de probabilidad, la probabilidad del subconjunto no puede exceder la probabilidad de ninguno de los superconjuntos que lo contienen:
		\begin{align*}
			F_{X,Y}(x, y) \le F_X(x) \quad \text{y} \quad F_{X,Y}(x, y) \le F_Y(y) \implies F_{X,Y}(x, y) \le \min\left( F_X(x), F_Y(y) \right).
		\end{align*}
		Para demostrar la cota inferior, invocamos la Desigualdad de Boole para la unión de eventos complementarios: la suma de las probabilidades de dos sucesos cualesquiera se relaciona con la probabilidad de su intersección mediante $P(A \cap B) = P(A) + P(B) - P(A \cup B)$.
		Haciendo $A = \set{X \le x}$ y $B = \set{Y \le y}$:
		\begin{align*}
			F_{X,Y}(x, y) = P(A \cap B) = F_X(x) + F_Y(y) - P(X \le x \text{ o } Y \le y).
		\end{align*}
		Como la medida de la unión de eventos en un espacio de probabilidad está axiomáticamente acotada por la unidad ($P(A \cup B) \le 1$), al restar un número menor o igual a $1$ se obtiene una cota inferior al sustituir por el valor extremo $1$:
		\begin{align*}
			F_{X,Y}(x, y) \ge F_X(x) + F_Y(y) - 1.
		\end{align*}
		Además, como por el primer axioma de Kolmogorov ninguna probabilidad acumulada puede ser estrictamente negativa ($F_{X,Y}(x, y) \ge 0$), combinamos ambas restricciones en un solo operador de máximo:
		\begin{align*}
			F_{X,Y}(x, y) \ge \max\left( 0, F_X(x) + F_Y(y) - 1 \right).
		\end{align*}
		En conclusión, la función de distribución acumulada conjunta discreta (y de cualquier tipo general) queda invariablemente confinada entre las dos cotas extremas universales de Fréchet-Hoeffding:
		\begin{align*}
			\max\left( 0, F_X(x) + F_Y(y) - 1 \right) \le F_{X,Y}(x, y) \le \min\left( F_X(x), F_Y(y) \right).
		\end{align*}
	\end{enumerate}
\end{solucion}
